\section{Reproducibility: code, records, deposit and provenance}
\label{si:repro}

\subsection{Code release and run records}

The code is released at \url{https://github.com/pa1aash/dft-gnn} under the MIT licence. Annotated tags mark the
states of the repository that the results depend on (Table~\ref{tab:tags}); the results reported in the article are
frozen at \tagresults{}, and the final code tag of the release is \texttt{[FINAL-TAG-TBD]}, to be set when the
deposit is minted.

\begin{table}[htbp]
\caption{Annotated tags of the code repository and what each marks (\texttt{git tag -n1}).}
\label{tab:tags}
\centering\small
% Generated by paper/si/tools/make_si_tables.py from the annotated tags of the repository (git tag -n1). Do not edit.
\begin{tabular}{@{}>{\raggedright}p{0.3\linewidth}>{\raggedright\arraybackslash}p{0.64\linewidth}@{}}
\toprule
Tag & What it marks \\
\midrule
\texttt{prereg-v1} & analysis plan frozen (registration of v1) \\
\texttt{tuned-v1} & tuned values of the originally registered network frozen \\
\texttt{g1-v1} & primary analysis of the learning curves (G1) \\
\texttt{nrp-code-v1} & pinned code of the cluster stages; the run ids of every cluster spec are computed with this name \\
\texttt{inputs-v1} & same commit as nrp-code-v1; release holding the processed inputs, checkpoints and MACE weights \\
\texttt{ingest-v1} & collaborator cluster analysis merged \\
\texttt{nrp-code-v2} & bug-fix tag of nrp-code-v1 for the relax, embed and geomeval reruns; run ids keep nrp-code-v1 \\
\texttt{nrp-code-v3} & result records also store the matgl version (record keeping only; no cluster run uses it) \\
\texttt{results-v1} & results freeze; headline numbers in results/freeze\_table.json \\
\texttt{prov-tune-v2, prov-prod-v2, prov-ckpt-diag} & code commits of the collaborator's v2 tuning, v2 production and checkpoint-diagnostic records \\
\bottomrule
\end{tabular}

\end{table}

Every training or analysis run writes one JSON record through a single function, \texttt{write\_result}, from a
clean working tree; the record holds the full configuration, the commit of the code, the package versions, the
seeds, the device (and, for the cluster runs, the GPU model) and the time, and the record of a training run names the hash of its per-site prediction table
(since \tagnrpthree{} it also stores the matgl version). A training run is named by its run id, a prefix of the
\shatwo{} digest of a canonical description of the run: model, hyperparameters, split, resample, budget, seed, code
reference and the hash of the graph-store manifest, plus any non-default option such as an epoch override, a
descriptor ablation or the component runs of the staged model (\texttt{src/dftgnn/train/\_\_init\_\_.py},
\texttt{run\_id}). The same description therefore always gives the same id, and two runs with the same id are the
same computation. For the cluster stages the code reference in the id is the tag \tagnrpone{} instead of the
generating commit, so the ids computed on the cluster and on the workstation agree; a worker runs such a job only if
the tag is an ancestor of its checkout and the code, configuration, split and job-specification directories are
unchanged since the tag. The bug-fix tag \tagnrptwo{} keeps the run ids of \tagnrpone{}. Named analysis records
(for example \texttt{results/g1\_primary.json}) are read by \texttt{paper/tools/make\_numbers.py}, which turns every
number of the article and of these Supplementary Materials into a macro with its source pointer; no number in
either document is typed by hand.

\subsection{Deposit and data availability}

The data and models are deposited in Zenodo under \ccby{} (DOI \texttt{[ZENODO-DOI-TBD]}, to be minted at
submission; \texttt{deposit/README.md}). The record holds the processed site table of the universe, the graph store
of the pristine supercells with its manifest, the release supercells of the hosts as released, the unit cells and
supercell tiling used for the machine-learned relaxations, all splits, the job specifications of the cluster stages,
the per-site prediction table and the record of every run, the checkpoints trained before the cluster stages, the
MACE-MP-0 weights used for the relaxations, and a \shatwo{} manifest of every file. Checkpoints, relaxed structures and
embeddings of the cluster runs are too large for the record; they are listed by \shatwo{} and are available from the
corresponding author on request. All formation energies, descriptors and structures derive from the oxygen-vacancy
database of Kumagai et al.\ \citep{Kumagai2021insights}, used at commit \releasecommit{} under \ccby{}; no new DFT
calculation was made.

\subsection{Provenance labels}

Each result group carries one of three labels, assigned by chronology rather than by outcome
(Table~\ref{tab:provenance}). The originally registered network (v1) and its rules were fixed in the analysis plan
at \tagprereg{} (\dateprereg{}) and its tuned values at \tagtuned{} (\datetunedvone{}). The primary network (v2) was
proposed after test-set diagnostics of v1 (design rule \datevtwoprotocol{}), selected on validation data only
(\datevtwoselect{}), retuned and frozen (\datevtwofreeze{}), and designated primary by author decision on
\dateauthordecision{}; every v2 result is post hoc and none is described as registered.

{\small
% Generated by paper/si/tools/make_si_tables.py from docs/provenance_labels.md. Do not edit.
\begin{longtable}{@{}r>{\raggedright}p{0.36\linewidth}>{\raggedright}p{0.2\linewidth}>{\raggedright\arraybackslash}p{0.3\linewidth}@{}}
\caption{Provenance label of every result group, transcribed from \texttt{docs/provenance\_labels.md} (groups 1 to 19: the collaborator analysis at its ingest; groups 20 to 35: the analyses added at the results freeze). PRE-REGISTERED: the originally registered backbone (v1) with the deciding rule committed before the results; POST HOC: the primary backbone (v2), or a rule committed after the results it decides; EXTENSION: compute or analysis beyond the registered plan. The labels record chronology, not quality. The records, rule commits and first-result times of every group are in the same file.}\label{tab:provenance}\\
\toprule
\# & Group & Network & Label \\
\midrule
\endfirsthead
\toprule
\# & Group & Network & Label \\
\midrule
\endhead
\bottomrule
\endlastfoot
1 & G1 primary (C1, C2), pre-registered arm & v1 & PRE-REGISTERED \\
2 & G1 on the v2 arm & v2 & POST HOC \\
3 & LOCO, S and D (\S{}11) & v1 (a654 values) & PRE-REGISTERED \\
4 & LOCO baselines B0 and RF-Kumagai on the folds & none & EXTENSION \\
5 & LOCO on the v2 arm & v2 & POST HOC \\
6 & staging (C3a): P1-v2 training and P-v2 evaluation & v2 (P1-v2 = \texttt{tuned-v1} P1 values + \texttt{init: kaiming}, not re-tuned) & POST HOC \\
7 & geometry evaluation (C3b) & v2 (S-v2, D-v2, P-v2 at B = 654 and 200) & POST HOC \\
8 & MLIP relaxations & none & PRE-REGISTERED \\
9 & latent probe (C4) & v2 & POST HOC \\
10 & screening (C5), pre-registered S & v1 (sweep S at B = 654) & PRE-REGISTERED \\
11 & screening (C5), S-v2 column & v2 & POST HOC \\
12 & descriptor ablations and permutation importances & v1 & EXTENSION \\
13 & D-late trained in the sweep protocol & v1 (D-late \texttt{tuned-v1} values) & EXTENSION \\
14 & budget $\times$ tuned-setting crossing & v1 & EXTENSION \\
15 & checkpoint gradient and flag diagnostics & v1 (retrained diagnostic checkpoints) & EXTENSION \\
16 & C0 tuning-overlap check & v1 (official C0 predictions) & EXTENSION \\
17 & v2 design screen & v2 candidates & POST HOC \\
18 & v2 retuning & v2 & POST HOC \\
19 & v2 production & v2 (\texttt{configs/\allowbreak{}tuned\_\allowbreak{}v2.yaml}) & POST HOC \\
20 & staging (C3a), registered analysis, step A2 & v1 & PRE-REGISTERED; the four-outcome TOST rule that decides its outcome is the collaborator rule of 2026-10-09, ratified 2026-10-11, POST HOC (as in \#6) \\
21 & staging (C3a-v2), re-analysis of \#6 with the G1 estimator, step A2 & v2 & POST HOC \\
22 & geometry evaluation (C3b) on our 818 relaxations at B = 654, with the non-converged sensitivity, step A3 & v1 & PRE-REGISTERED \\
23 & geometry at B = 200 & v1 & DESCRIPTIVE EXTRA, not pre-registered (label of \texttt{docs/\allowbreak{}a3\_\allowbreak{}notes.md}, kept) \\
24 & geometry: TOST of $\Delta$MAE\_P - $\Delta$MAE\_S and slope differences & v1 & POST HOC \\
25 & geometry: cross-check against the collaborator's 809 relaxations and his v2 evaluation & his v2 (\texttt{results/\allowbreak{}mlip\_\allowbreak{}eval\_\allowbreak{}v2.json}, \#7) & POST HOC (descriptive) \\
26 & latent probe (C4), registered analysis, step A4 & v1 (60 seed-0 S checkpoints of the sweep) & PRE-REGISTERED \\
27 & probe: within-host R$^{2}$ and per-descriptor values (PRB-WH) & v1 & DESCRIPTIVE (not part of the C4 rule; label of \texttt{docs/\allowbreak{}claims.yaml}, kept) \\
28 & probe on S-v2 (C4-v2), recomputed from \#9 & v2 & POST HOC (descriptive) \\
29 & precision extension (N4), step A1 & v1 & EXTENSION \\
30 & precision extension (N4), v2 & v2 & POST HOC (v2 backbone); EXTENSION in compute \\
31 & grouped K-fold CV (N3), step A1 & v1 & EXTENSION \\
32 & grouped K-fold CV (N3), v2 & v2 & POST HOC (v2 backbone); EXTENSION in compute \\
33 & high-moment sensitivity (SENS1), step A5 & v1 & EXTENSION \\
34 & CGCNN cross-check (SENS2), untuned, descriptive, step A5 & its own architecture; MEGNet v1 and v2 beside it & EXTENSION \\
35 & capped 600-epoch twins (SENS3), step A5 & v1 & POST HOC \\
\end{longtable}
}

\subsection{Deviations from the registered plan}

Every departure from the analysis plan is logged in \texttt{docs/deviations.md} before the run or analysis it
affects, with the plan section, the reason and the impact, and no row is edited after it is written; results
generated before a deviation are kept and reported. The rows fall into six kinds. Clarifications fix what the plan
left open without changing a registered decision: the early-stopping threshold, the seed rules, the tuning protocol,
the bootstrap and the rule for the smallest qualifying budget, the definitions of the geometry, probe and
leave-one-chemistry-out analyses, and operational details of the job queue. Deviations add an arm or an analysis
without removing one: the unselected injection route as a sensitivity, and the v2 network with its design screen,
selection rule and use in the staging, geometry and probe analyses. Extensions add compute or analysis beyond the
plan, each logged before its first result: the precision resamples, the grouped cross-validation, the within-host
breakdown, the rerun of the relaxations with an extended tiling rule, the repeatability study and the analyses of
the results freeze. One correction restricts the statement that the training cap never binds to the registered
check of the cap, because runs of the registered sweep did stop at the cap. One post hoc ratification records the
author sign-off, missing from the repository, of three collaborator rows: the equivalence rule and four-outcome
classification of the staging and geometry comparisons, the reporting rules for the smallest qualifying budget, and
the use of v2 for the staging, geometry and probe analyses; all three stay post hoc. Two author decisions designate v2 as the primary network and freeze the results at \tagresults{}; after the
freeze no experiment, run or analysis was added. The full log, with the commit time of every row, is part of the
code release.
